% !PRINTED 221
\weberpage{221}

Die Berechnung der Co\"efficienten $Q$ ist leicht auszuführen, wenn man diese Gleichung für $x_1$, $x_2$, $x_3$ und entsprechend $y_1$, $y_2$, $y_3$ aufstellt.

Wir wollen über den Gang der Rechnung, die nur die Darstellung symmetrischer Functionen der Wurzeln einer cubischen Gleichung durch die Co\"efficienten nach den Vorschriften des vierten Abschnittes (\S.\,42, 45) erfordert, einige Andeutungen machen. Nimmt man zunächst die Summe der drei Gleichungen (9), so erhält man nach (6)
\[
3 Q_0 = \frac{2}{3} H(t_1, t_0) Q_2 - a_1 f(t_1, t_0),
\]
so dass also nur noch $Q_1$ und $Q_2$ berechnet zu werden brauchen; diese findet man, wenn man die für $x = x_1$ gebildete Gleichung (9) mit
\[
y_2 - y_3 = a_0 (x_2 - x_3) (t_1 - t_0 x_1)
\]
und mit
\[
y_2^2 - y_3^2 = -a_0 (x_2 - x_3) y_1 (t_1 - t_0 x_1)
\]
multiplicirt und die drei durch cyklische Vertauschung der Indices 1, 2, 3 gebildeten analogen Gleichungen addirt.

Man hat dann nur Gebrauch zu machen von den beiden Formeln
\[
a_0^2 \Sigma x_1^2 (x_2 - x_3) = -\sqrt{D}
\]
\[
a_0^3 \Sigma x_1^3 (x_2 - x_3) = a_1 \sqrt{D},
\]
und findet so
\[
Q_0 = \frac{-a_1}{3} f(t_1, t_0) - \frac{2}{3} t_0 H(t_1, t_0)
\]
\[
Q_1 = a_0 t_1^2 + \frac{2}{3} a_1 t_0 t_1 + \frac{1}{3} a_2 t_0^2
\]
\[
Q_2 = -t_0.
\]

Setzen wir zur Vereinfachung $t_1 = t$, $t_0 = 1$, gehen also zu den inhomogenen Ausdrücken über, so ergiebt sich
\weq{10}{
(3 a_0 x + a_1) f(t) = -\frac{2}{3} H(t) + y f'(t) - 3 y^2,
}
während die Gleichung für $y$ lautet
\weq{11}{
y^3 + \frac{1}{3} H(t) y + \frac{1}{27} Q(t) = 0,
}
wenn wir $H(t)$, $Q(t)$ für $H(t, 1)$, $Q(t, 1)$ setzen.

Um nun die cubische Gleichung zu lösen, bestimmen wir $t$ aus der quadratischen Gleichung
\[
H(t) = 0.
\]
